% Criterion 5 chapter converted from the approved Word source.
% Expected packages in the main document: array, booktabs, longtable,
% multirow, graphicx, hyperref, calc, enumitem, and float.
\providecommand{\real}[1]{#1}
%\usepackage{multirow}


\hypertarget{criterion-5-faculty-and-support-staff}{%
\chapter{Criterion 5: Faculty and Support Staff}\label{criterion-5-faculty-and-support-staff}}

The Department of Software Engineering at SEECS, NUST, maintains a complement of dedicated and shared faculty, supported by lab engineers, to facilitate effective student-faculty interaction and the delivery of the undergraduate program curriculum. Faculty members possess diverse academic qualifications, professional experience, and areas of specialization to support the core and interdisciplinary components of the Software Engineering curriculum. This criterion presents the department's faculty strength, professional ranks, academic qualifications, areas of expertise, teaching responsibilities, and student-faculty ratio, together with the supporting staff resources required for effective program delivery. The relevant faculty duties and workload arrangements are discussed in the subsequent sections, with reference to the applicable NUST statutes and policies (Annexure~\ref{annex:AG}).

\section{Faculty Strength and Competencies}
\label{faculty-strength-and-competencies}

The department currently has 18 full-time dedicated faculty members, comprising one Professor, two Associate Professors, 12 Assistant Professors, and three Lecturers (Table~\ref{tab:headcount1}). In addition, six shared faculty members contribute to teaching in their respective areas of expertise, while five Lab Engineers support practical instruction and laboratory activities. The faculty complement provides coverage across Software Engineering, Artificial Intelligence, Information Security, and Applied Computing, supporting the core components of the program. Shared and visiting faculty from other disciplines support the interdisciplinary components of the program. The faculty's academic qualifications, professional designations, institutional backgrounds, experience, departmental affiliations, and joining dates are documented in Annexure~\ref{annex:I}. Annexure~\ref{annex:K} provides the current and preceding semesters' teaching workloads, distinguishing theory and practical credit hours for undergraduate and postgraduate courses, where applicable. %The list of faculty members registered with the Pakistan Engineering Council (PEC) is provided in Annexure~\ref{annex:M}.

\begin{table}[htbp]
\centering
\small
\setlength{\tabcolsep}{4pt}
\caption{Faculty Headcount by Category}
\label{tab:headcount1}
\begin{tabular}{@{}p{0.34\columnwidth}ccccc@{}}
\toprule
\textbf{Category} & \textbf{Prof.} & \textbf{Assoc. Prof.} &
\textbf{Asst. Prof.} & \textbf{Lect./LE} & \textbf{Total} \\
\midrule
Full-time dedicated faculty & 1 & 2 & 12 & 3 & 18 \\
Shared faculty              & 1 & 1 & 4 & -- & 6 \\
Visiting Engg. faculty              & -- & 1 & -- & -- & 1 \\
Lab Engineers               & -- & -- & -- & 5 & 5 \\
\bottomrule
\end{tabular}
\end{table}

\subsection{Faculty Qualifications and Professional Ranks}
\label{pyramid-of-academic-structure}

The department maintains a mix of academic ranks and qualifications to support curriculum delivery, student supervision, and academic development. Currently, there are 18 FTDF (Engineering and non-engineering) members, comprising one Professor, two Associate Professors, 12 Assistant Professors, and three Lecturers. As reported in the faculty qualification records, 11 FTDF members, hold PhD degrees, while seven FTFD members hold MS qualifications. The institute provides opportunities for professional advancement based on applicable policies and eligibility criteria. Three of the seven MS FTDF members are concurrently pursuing PhD degrees (Table~\ref{tab:high}). One of the lab engineers is also pursuing PhD degree. 

\begin{longtable}{|>{\centering\arraybackslash}p{1cm}
|>{\raggedright\arraybackslash}p{2.8cm}
|>{\raggedright\arraybackslash}p{1.7cm}
|>{\raggedright\arraybackslash}p{2.5cm}
|>{\raggedright\arraybackslash}p{3cm}|}

\caption{SE Faculty Pursuing PhD}
\label{tab:high}\\
\hline
\textbf{Sr. No.} &
\textbf{Name} &
\textbf{Host Country} &
\textbf{University} &
\textbf{Subject} \\
\hline
\endfirsthead

\hline
\textbf{Sr. No.} &
\textbf{Name} &
\textbf{Host Country} &
\textbf{University} &
\textbf{Subject} \\
\hline
\endhead

1 & Hirra Anwar & Pakistan & SMME-NUST & Robotics and AI \\
\hline

2 & Mobeena Shahzad & Pakistan & EME-NUST & Robotics and AI \\
\hline

3 & Ayesha Kanwal & Pakistan & SMME-NUST & Robotics and AI \\
\hline

\end{longtable}

\subsection{Full-Time Dedicated Faculty}
\label{full-time-dedicated-faculty}

The 18 full-time dedicated faculty members constitute the department's primary teaching and academic resource. Their qualifications, experience, and areas of expertise support the delivery of the Software Engineering curriculum, student advising, project supervision, and related academic activities. Their teaching contributions and workload distribution across the current and preceding semesters are documented in Annexure~\ref{annex:K}.

\subsection{Shared Faculty}
\label{shared-faculty}

Six non-engineering shared faculty members supplement the dedicated faculty by teaching courses within their respective areas of specialization. These faculty members are associated with other departments or academic programs and enable the department to draw on expertise beyond its dedicated faculty complement. Their designations, qualifications, experience, and departmental affiliations are recorded in Annexure~\ref{annex:I}, while their teaching contributions are documented in Annexure~\ref{annex:K}.

\subsection{Visiting and Adjunct Faculty}
\label{visiting-faculty}

The department engages Temporary Visiting Faculty (TVF), where required, to address specialized teaching needs that cannot be fully met through the available permanent faculty. Such appointments are particularly relevant to elective courses requiring industry-specific knowledge and professional experience. In appropriate cases, NUST PhD scholars may also be engaged as TVFs. University policy further permits eligible permanent faculty members to undertake additional teaching assignments as visiting faculty alongside their regular responsibilities. These arrangements provide flexibility in addressing specialized curricular requirements while maintaining the department's core reliance on its dedicated and shared faculty. There are two adjunct faculty members in the department as well namely Dr. Elena Gaura from Coventry University, UK, and Dr. Ali Imran Shariq from the Norwegian University of Science and Technology. 

\subsection{Student--Faculty Ratio}
\label{student-teacher-ratio}

The department has maintained a healthy student--faculty ratio since the undergraduate program commenced in 2010 (Annexure~\ref{annex:J}). The current student enrolment is 349, supported by 18 full-time dedicated faculty members, six shared faculty members, and one TVF. Five Lab Engineers provide additional support for laboratory-based teaching and practical activities. As summarized in Table~\ref{tab:ratio}, the student--faculty ratio is 15.44:1 when calculated against current enrolment and the faculty headcount (22.6) used for this calculation. Against the reported allowed student intake, the ratio is 15.92:1. Both ratios are better than the PEC guideline of 20:1, demonstrating that the reported faculty strength is sufficient to support effective student--faculty interaction. These ratios are reported in accordance with the applicable PEC accreditation criteria.

\begin{table}[!h]
\centering
\small
\setlength{\tabcolsep}{4pt}
\caption{Student-Faculty Ratio}
\label{tab:ratio}
\begin{tabular}{@{}p{0.34\columnwidth}ccccc@{}}
\toprule
\textbf{Category} & \textbf{PhD} & \textbf{MS/MSc} &
\textbf{BE/BSc} & \textbf{Total} & \textbf{PEC Criteria} \\
\midrule
Full-time dedicated faculty & 9 & 6 & 0 & 15 &  \\
(Engineering)&&&&&\\
\hline
Full-time dedicated faculty & 2 & 1 & 0 & 1.6 & Capped at 20\% of \\
(Non-Engineering)&&&&&MFS (8 FTDF/Sect.)\\
\hline
Shared faculty              & 4 & 2 & 0 & & Counted half \\
(Non-Engineering)&&&&3.5&Capped at\\
Shared faculty              & 1 & 0 & 0 & & 25\% of FTDF \\
(Engineering)&&&&&\\
\hline
Lab Engineers               & 0 & 3 & 2 & 2.5 & Counted half \\
&&&&&Capped at 20\%\\
&&&&&FTDF\\
\hline
\textbf{Total Faculty Count}&\textbf{16}&\textbf{12}&\textbf{2}&\textbf{22.6}&\\
\midrule
\textbf{Allowed Intake}&\multicolumn{5}{c}{360} \\
\textbf{Current Student Strength}&\multicolumn{5}{c}{349} \\
\midrule
Student--faculty ratio & \multicolumn{5}{c}{15.44:1} \\
(As per current enrollment)& \multicolumn{5}{c} {}\\
Student--faculty ratio & \multicolumn{5}{c}{15.92:1} \\
(As per allowed enrollment)&\multicolumn{5}{c} {}\\
\bottomrule
\end{tabular}
\end{table}

Overall, the department's dedicated and shared faculty resources, supported by laboratory personnel and specialist visiting appointments where needed, provide a basis for effective curriculum delivery and student engagement. The documented faculty qualifications, professional ranks, expertise, and teaching workloads provide supporting evidence of the department's capacity to cover the program's curricular requirements.

%All members of the NUST fraternity assigned to teaching, research, and academic management activities comprise faculty. The faculty at SEECS includes the Principal, Associate Deans, Heads of Departments, Professors, Associate and Assistant Professors, and Lecturers. The faculty duties have been specified in NUST Statutes VI and are attached as Annexure~\ref{annex:AG}. 


%Faculty positions since the start of the undergraduate program in 2010 have been improving gradually. Relevant data since the last visit of the PEC Accreditation Committee, related to faculty summary, is given in Annex~\ref{annex:J}. Total enrolled students presently are 348. There are 14 full-time dedicated engineering faculty members. Full-time dedicated non-engineering faculty members are 3. Based on the criteria given in section 3.2.5.3 of the revised PEC accreditation manual, the student-teacher ratio comes to 17.53. Detailed calculations are given in Table 6.1 below.



\section{Training and Mentoring on Pedagogical Skills and OBE}
\label{sec:training}

SEECS has recently formalized a staged mechanism for training and mentoring faculty members in pedagogical skills and Outcome-Based Education (OBE) concepts. Formalized through the SEECS Faculty On-Boarding Guide (Version 1.0, published September 11, 2026\footnote{Faculty-On-Boarding-Guide-v1.0-1.pdf}), this framework operates at the Department and School level (provided in Evidence Folder). Newly inducted faculty members must also complete a NUST-level on-boarding training prior to and during their first semester of teaching conducted by NUST Center for Teaching and Learning (NCTL).

\subsection{Multi-Tier Mentoring and Training Mechanism}
Faculty training operates across a multi-tiered structure:
\begin {itemize}
\item \textbf{NUST Level:} Newly appointed faculty undergo a structured onboarding program covering    instructional delivery, class conduct, effective communication techniques, teaching methodologies, and lecture preparation. Operational training is conducted for the NUST LMS, Qalam CMS, and E-Office platforms. These mandatory sessions are organized by NCTL.

\item \textbf{(SEECS) Level:} At the school level, SEECS facilitates continuous training and guidance through:

\begin{itemize}
    \item \textbf{Knowledge Groups:} Knowledge Group Heads (KGHs) provide domain-specific oversight, vetting course outlines, mid-semester exams (MSE), end-semester exams (ESE), and retake papers for alignment with OBE standards.
    \item \textbf{Mentoring:} A senior faculty mentor is assigned to guide the inductee throughout their first semester in syllabus delivery, lab conduct, assessment vetting, and grading criteria.  
\end {itemize}
    
\end{itemize}

\subsection{OBE Concepts and Implementation Methodologies}

Faculty follow a systematic mechanism to gain a thorough operational understanding of Program Educational Objectives (PEOs), Program Learning Outcomes (PLOs), and Course Learning Outcomes (CLOs). 
\begin {itemize}[noitemsep, topsep=0pt]
\item \textbf{OBE Trainings:} OBE trainings are frequently held to enhance the understanding of faculty members regarding OBE framework followed at SEECS (Annexure~\ref{annex:T}). Each faculty member provides an overview of OBE framework to the students of their respective classes at the beginning of the semester. Students are also detailed during orientation session.
\item \textbf{Assessment Mapping:} Faculty builds a detailed assessment map linking each evaluation item (quizzes, assignments, exams, and practical labs) directly to approved CLOs and Bloom's difficulty levels mentioned in the approved course outline. 
\item \textbf{System Integration:} OBE data is configured in the Qalam Campus Management System (CMS) gradebook to aggregate mapped marks and generate real-time direct CLO–PLO attainment reports. 
\item \textbf{Continuous Quality Improvement (CQI):} Course folders maintained on the Learning Management System (LMS) include attainment reports and instructor CQI action plan to address any outcomes falling below target thresholds.
\end{itemize}

\subsection{Effectiveness of Training and Mentoring}
The efficacy of the training framework is measured through multi-faceted feedback loops:   
\begin{itemize}[noitemsep, topsep=0pt]
    \item \textbf{Student Course Evaluations:} Conducted digitally via Qalam at the end of each semester to track student satisfaction and instructional quality over successive terms.
    \item \textbf{Peer Teaching Audits:} Periodic peer reviews assess lecture clarity, responsiveness, classroom management, and alignment with approved outlines.
    \item \textbf{CLO/PLO Attainment Metrics:} Automated Qalam reports evaluate cohort performance per outcome, triggering mandatory CQI interventions where target thresholds are missed. 
\end{itemize}

\section{Faculty Development for Professional Growth}
Faculty's professional growth is supported through institutional faculty development, professional training, mentoring, research engagement, and opportunities for academic and professional advancement. Faculty development activities include pedagogical and Outcome-Based Education (OBE) training, technical and research-oriented workshops, participation in national and international professional development programmes, and guidance provided through departmental and institutional mechanisms. 

NUST supports the professional growth of its faculty through the Professional Development Centre (PDC), established in July 2007 under an initiative of the Higher Education Commission (HEC). Recently, it has been developed into Iqbal Executive Development Centre (IEDC) that offers professional training using experienced facilitators, practical exercises, and case-based learning to strengthen professional competencies and institutional practices. Faculty development activities provide opportunities to enhance pedagogical skills, OBE implementation, research capabilities, and academic leadership. The effectiveness of these activities can be assessed through their application in teaching and research, career progression, faculty feedback, and retention trends, enabling the identification of further development needs and appropriate improvement measures. Evidence of these will be provided separately. Table~\ref{tab:training} shows some recent professional development trainings in which SE faculty participated.

\begin{longtable}{|>{\centering\arraybackslash}p{1cm}
|>{\raggedright\arraybackslash}p{4cm}
|>{\centering\arraybackslash}p{2cm}
|>{\raggedright\arraybackslash}p{2.7cm}
|>{\raggedright\arraybackslash}p{3.5cm}|}
\caption{Training and development of SEECS faculty}
\label{tab:training}
\hline
\textbf{Sr. No} & 
\textbf{Training/Short Course} & 
\textbf{Date} & 
\textbf{Attendee(s)} & 
\textbf{Organized by} \\ 
\hline
\endhead

1 & 
AI in Agriculture & 
3/23/2026 & 
Dr. Momina Moetesum & 
Asian Productivity Organization \\ 
\hline

2 & 
Reinforcement Learning Unleashed: Hands-On Skill Building for Smarter AI & 
08/18/2025 to 08/22/2025 & 
Dr. Zuhair Zafar & 
Machine Vision and Intelligent Systems (MachVIS) Lab, SEECS, National University of Sciences and Technology, Islamabad. \\ 
\hline

3 & 
AI for Manufacturing Industry & 
4/23/2025 & 
Dr. Momina Moetesum & 
Continuing Professional Development Program, Department of Electrical Engineering, Lahore College for Women University, Lahore. \\ 
\hline

4 & 
IEEE Learning \& Research for Pakistan students \& researchers: an overview of IEEE Xplore Digital Library & 
4/15/2026 & 
Dr. Momina Moetesum & 
A\&J Def Tek Consultants \\ 
\hline

5 & 
National Workshop on AI Use-case Development & 
11/12/2025 & 
Dr. Momina Moetesum & 
IGNITE \\ 
\hline

6 & 
"Research Smarter: Using Web of Science for Academic Success" & 
5/21/2026 & 
Dr. Faisal Shafait & 
Higher Education Commission (HEC) Pakistan \\ 
\hline

7 & 
IEEE Authorship and Open Access Symposium: Tips and Best Practices to Get Published from IEEE Editors & 
6/12/2026 & 
Dr. Ayesha Maqbool & 
IEEE \\ 
\hline

\end{longtable}

\hypertarget{faculty-workload}{%
\section{Faculty Workload Distribution}\label{faculty-workload}}
Faculty workload and general duties are covered in detail in Annexure~\ref{annex:AG}. All Faculty members, besides teaching, are also involved in other activities which primarily include:

\begin{itemize}[noitemsep, topsep=0pt]
\item Performing administrative duties 
\item Performing the duties of LQEC
\item UG FYDP Supervision 
\item MS \& PhD Thesis Supervision
\item Staff appointments
\item Program coordination
\item Invigilation 
\item Funded projects
\item Research publications
\end{itemize}

In most cases, regular teaching load is reduced to facilitate these activities. For instance, a waiver of 1 credit hour is awarded on supervision of a PhD scholar. Similarly, administrative tasks like heading a department is equivalent to 3 credit hours. Nonetheless, maximum non-teaching credit is equivalent to 6 credit hours. To streamline the equivalence formula for equating teaching load with administrative duties, NUST Academic policy is relied upon using which a faculty member's research and administrative duties are elucidated by assigning them credit hours. The total load of a faculty member is the sum of Teaching, Research and Administrative duties. 

\hypertarget{faculty-retention-and-career-planning}{%
\section{Faculty Retention and Career Planning}\label{faculty-retention-and-career-planning}}

NUST provides remuneration and benefits in accordance with the applicable
institutional statutes, service regulations, and employment terms. The NUST
Gazette and Statutes are available through the University's official
website.\footnote{NUST, \textit{NUST Gazette and Statutes}, available at:
\url{https://nust.edu.pk/downloads/nust-gazette-and-statutes/}.}
Faculty policies are summarised in the NUST Faculty Handbook\footnote{NUST,
\textit{Faculty Handbook}, Version 1.3, Chapters 3, 5, 6 and 8. Detailed rules
are contained in the NUST Human Resource Handbook.} and administered by the
Human Resource Directorate. The applicable pay arrangements depend on the
faculty member's employment category:

\begin{enumerate}
    \renewcommand{\labelenumi}{(\alph{enumi})}
    \item \textbf{NUST Faculty:} Salaries and benefits are determined in
    accordance with the applicable NUST Statutes and institutional
    regulations. Faculty may be appointed on the Tenure Track System (TTS),
    which NUST has adopted in parallel with the HEC framework since 2012.

    \item \textbf{Ministry of Defence Officers:} Remuneration is administered
    in accordance with the applicable Government of Pakistan pay scales and
    regulations.

    \item \textbf{Pakistan Army Officers:} Pay is administered according to
    the applicable military pay scales through the Controller Military
    Accounts.
\end{enumerate}

In addition to remuneration, NUST provides the following facilities and
benefits to support faculty welfare, professional development, and retention.

\begin{enumerate}
    \item \textbf{Career Progression and Performance Appraisal.}
    Faculty performance is reviewed annually through the Annual Appraisal
    Report (AAR), which combines self-assessment, student feedback collected
    twice per semester, class audit, and evaluation by the Head of Department
    and Dean. Research performance (publications, sponsored funding,
    commercialisation, and student supervision) is a weighted input to
    subsequent promotion decisions.

    \item \textbf{Leave.}
    Faculty are entitled to 20 days of casual leave and 48 days of earned
    leave per year. Leave for study or academic purposes is approved by the
    Pro-Rector (Academics).

    \item \textbf{On-Campus Accommodation and Security.}
    Residential accommodation is available to faculty members on a
    first-come, first-served basis in the Iqra and Isra apartments and villas
    of various sizes, as well as faculty messes for single occupants. The
    campus is protected by CCTV surveillance and guards at all entries and
    exits.

    \item \textbf{Medical Facilities and Health Insurance.}
    NUST operates an in-house Medical Center providing consultation,
    emergency treatment, ambulance service, and a pharmacy. In addition,
    contractual civilian faculty members are covered, together with their
    eligible dependents (spouse and children), under the NUST Group Health
    Insurance Scheme.\footnote{NUST, \textit{NUST Group Health Insurance
    Scheme for Contractual Faculty and Support Staff Officers (Civilian
    only): General Guidelines/Instructions}, Human Resource Directorate.}
    Enrolment and records are managed by the Human Resource Directorate, the
    NUST Medical Centre acts as the liaison office with the insurance company
    for card issuance, activation and reimbursement claims, and the Finance
    Directorate processes premium payments. Dependent eligibility lapses on
    the death of the dependent, on separation (spouse), on marriage (son or
    daughter), or when a son exceeds 23 years of age. Coverage is suspended
    during leave exceeding four months and resumes on return to duty. Where
    the insurance limit is exhausted, further medical expenditure may be
    claimed through the Medical Reserve Pool, subject to rationalisation at
    the rates prescribed by Armed Forces Hospitals. Medical benefits for
    faculty in other employment categories are provided in accordance with
    the applicable institutional policy and eligibility
    conditions.\footnote{NUST, \textit{Medical Policy 2015 (Revised)},
    available at:
    \url{https://nshs.nust.edu.pk/wp-content/uploads/2023/08/Medical-Policy.pdf}.}

    \item \textbf{Research and Publication Incentives.}
    NUST sponsors article publication charges for peer-reviewed journals as per policy\footnote{NUST Research Directorate, \textit{Publication Charges for
    Peer-Reviewed Journals Sponsored by Research Dte}, Working Paper
    No.~39, 61st ACM, 2021; and policy letters No.~0986/41/Research/NUST
    dated 16 February 2024 and 15 November 2024.} Only journals indexed in
    the Science Citation Index Expanded, Social Sciences Citation Index, or
    Arts and Humanities Citation Index of Web of Science and ranked in
    Quartile~1 or Quartile~2 of the Journal Citation Reports are eligible.
    Emerging Sources Citation Index journals are excluded. NUST may sponsor
    up to US\$1,800 for Quartile~1 articles and up to US\$1,200 for
    Quartile~2 articles, shared among the NUST authors by author order
    (sole author 100\%; two authors 60\%/40\%; three authors
    50\%/35\%/15\%; four authors 45\%/30\%/15\%/10\%). The quartile is taken
    from the JCR edition of the year of publication. Applications are
    submitted at least four weeks before the invoice due date and within
    three months of acceptance, with a Turnitin report and an undertaking
    against plagiarism, endorsed by the Principal of the institution. For
    Quartile~2 articles one case per faculty member is permitted per fiscal
    year. Invoices must be issued in the name of the NUST authors, with split
    invoices where foreign co-authors are involved. Exceptional cases may be
    considered for full charges. Faculty who claim publication charges are
    not eligible for a publication financial award. In addition, NUST
    sponsors registration fees and travel for peer-reviewed international
    conferences, subject to eligibility and availability of funds, and funds
    professional society memberships and journal subscriptions. Annual Best
    Teacher and Best Researcher awards carry prizes of Rs~75,000 at school
    level and Rs~150,000 at university level. Seed funding for research is
    available from the NUST R\&D Fund, and faculty are supported in obtaining
    external grants through the Research, Innovation and Commercialization
    Directorate.

    \item \textbf{Provident Fund.}
    Provident fund benefits are subject to the applicable institutional
    policies, employment terms, and eligibility requirements, which are
    maintained by the NUST Human Resource Directorate.

    \item \textbf{Service Gratuity.}
    NUST provides service gratuity in accordance with the applicable policy
    and eligibility requirements. The qualifying service period varies
    according to the employment category.\footnote{NUST, \textit{NUST
    Gratuity Policy}, Policy Code 08-00-001-04, available at:
    \url{https://nshs.nust.edu.pk/wp-content/uploads/2023/08/Gratuity.pdf}.}

    \item \textbf{Faculty Development.}
    The NUST Faculty Development Program, launched in 2002--03, sponsors
    faculty for PhD and MS-leading-to-PhD studies at reputed foreign
    universities. The 2025 prospectus reports that 430 scholars have been
    sent abroad, of whom 316 have rejoined. PhD-qualified faculty may be
    sponsored for 9--12 months of post-doctoral research abroad, subject to
    availability of funds. A structured continuing professional development
    programme of seven modules covers orientation, instructional
    development, research writing, administrative and personal development,
    and short courses.

    \item \textbf{Reasonable Faculty Workload and Support for Higher
    Studies.}
    The average teaching load of a faculty member does not exceed 12 credit hours and may be reduced for administrative responsibilities, approved research, laboratory and course
    development, and supervision of theses and projects. SEECS endeavours to
    distribute teaching and other academic responsibilities reasonably, in
    accordance with applicable PEC guidelines and institutional requirements.
    Where feasible, workload allocation also takes account of the academic
    development needs of early-career faculty members pursuing higher
    studies, while maintaining the continuity and quality of teaching,
    research, and departmental activities. Compliance with the applicable
    workload requirements is supported by departmental workload records.

    \item \textbf{Welfare and Family Support.}
    Counselling services through the Center for Counseling and Career
    Advisory (C3A) are free of charge for NUST faculty and their immediate
    families. The NUST Creative Learning School provides day care and
    schooling from pre-school to middle school, with after-care services at no
    extra charge. A campus and mass-transit transport system serves faculty
    from all localities of the twin cities, and partner banks offer employee
    packages including free ATM transactions.
\end{enumerate}

\section{Technical and Administrative Support Staff}

The SE programme is supported by technical and administrative staff
whose numbers and qualifications are adequate to deliver the programme. Programme-specific laboratory support is provided by five PEC-registered Lab Engineers, who are counted towards the student--teacher ratio at one-half each
in accordance with PEC guidelines. In addition,
SE students use the computing facilities, each of which has an assigned lab technician. Technicians are responsible for system installation,
maintenance, software deployment and lab-session support, which keeps the
laboratories up to date and resolves technical issues promptly. Although
lab technicians are not counted towards the student--teacher ratio, their
role is essential to uninterrupted laboratory operation. Administrative
support is provided by the SEECS Administration Department, Exam Branch,
Academic Coordination Branch (ACB), ITS Department, and Qalam/LMS support staff. Support staff
performance is reviewed through annual appraisal mechanism and are provided opportunities for professional development through courses and trainings.

\section{Faculty Research and Grants}
\label{researchfunding}
Faculty members at SEECS including SE, actively pursue research across diverse areas of computing and engineering, resulting in scholarly publications (Annexure~\ref{annex:r}) and research collaborations with national and international partners. Faculty research is further supported through sponsored projects funded by industry, government bodies, and other donor agencies, promoting knowledge creation, innovation, and the practical application of research outcomes. Evidence of the continuation of faculty research activities, publications, and sponsored projects is available in the Faculty Research Evidence folder. 

